\section{Divisor Gram positivity and the terminal-octave floor}
\label{sec:gram}

For a finitely supported sequence $u=(u_d)$, define
\begin{equation}
 \cQ(u)=\sum_{d,e}\frac{u_d\overline{u_e}}{[d,e]}.
 \label{eq:Gram-Q}
\end{equation}

\begin{lemma}[Selberg divisor Gram factorization]
\label{lem:Gram-factorization}
For every finitely supported $u$,
\begin{equation}
 \boxed{\quad
 \cQ(u)=\sum_{r\ge1}\varphi(r)
 \left|\sum_{r\mid d}\frac{u_d}{d}\right|^2\ge0.
 \quad}
 \label{eq:Gram-factorization}
\end{equation}
\end{lemma}

\begin{proof}
Use
\[
 \frac1{[d,e]}=\frac{(d,e)}{de}
 =\frac1{de}\sum_{r\mid(d,e)}\varphi(r)
\]
and interchange the finite sums.
\end{proof}

Take $u_d=\mu(d)\one_{d\in I}$ and write
\begin{equation}
 Q_I=\sum_{d,e\in I}\frac{\mu(d)\mu(e)}{[d,e]}.
 \label{eq:QI}
\end{equation}

\begin{theorem}[Terminal-octave rigidity]
\label{thm:terminal-octave}
If $D_0<V/2$, then
\begin{align}
 Q_I
 &\ge\sum_{V/2<r\le V}\frac{\mu(r)^2\varphi(r)}{r^2}
 \label{eq:terminal-octave-finite}\\
 &=C_{\mathrm{sf}}\log2+o(1),
 \qquad
 C_{\mathrm{sf}}=\prod_p
 \left(1-\frac2{p^2}+\frac1{p^3}\right)>0.
 \label{eq:terminal-octave-asymptotic}
\end{align}
\end{theorem}

\begin{proof}
In \eqref{eq:Gram-factorization}, retain only $V/2<r\le V$.
The only multiple of such an $r$ not exceeding $V$ is $d=r$; because
$D_0<V/2$, it belongs to $I$.  This proves
\eqref{eq:terminal-octave-finite}.

For completeness, put $a(n)=\mu(n)^2\varphi(n)/n$.  Its mean value is
\[
 \prod_p(1-p^{-1})\left(1+\frac{1-p^{-1}}p\right)
 =\prod_p\left(1-2p^{-2}+p^{-3}\right)=C_{\mathrm{sf}}.
\]
Standard mean-value theory for multiplicative functions, followed by
partial summation, gives
\[
 \sum_{y<r\le z}\frac{a(r)}r
 =C_{\mathrm{sf}}\log(z/y)+o(1)
\]
when $z/y$ is fixed and $y\to\infty$.  Take $y=V/2$, $z=V$.
\end{proof}

The finite-$K$ meaning of this Gram mass is exact up to an elementary
lattice error.

\begin{proposition}[Weighted second moment]
\label{prop:weighted-second-moment}
Let $w\in C_c^1((0,\infty))$.  Then
\begin{equation}
 \sum_{k\ge1}\beta_I(k)^2w(k/K)
 =K I_0(w)Q_I+O_w(V^2),
 \qquad I_0(w)=\int_0^\infty w(t)\,dt.
 \label{eq:weighted-beta-second-moment}
\end{equation}
In particular, if $w\ge0$, $I_0(w)>0$, $D_0<V/2$, and
$K/V^2\to\infty$, then
\begin{equation}
 \sum_k\beta_I(k)^2w(k/K)\gg_w K.
 \label{eq:beta-natural-L2}
\end{equation}
\end{proposition}

\begin{proof}
Insert \eqref{eq:lcm-transform} and use
\[
 \sum_jw(qj/K)=\frac Kq I_0(w)+O_w(1).
\]
The main coefficient is
\[
 \sum_q\frac{\Gamma_I(q)}q
 =\sum_{d,e\in I}\frac{\mu(d)\mu(e)}{[d,e]}=Q_I,
\]
while \eqref{eq:Gamma-l1} bounds the total error by $O_w(V^2)$.
The final assertion follows from \cref{thm:terminal-octave}.
\end{proof}

\begin{remark}[A rigorous coefficient-shape obstruction]
Suppose $D_0<\sqrt V$ and let $p\in(\sqrt V,V]$ be prime.  Then
$\mu(p)\one_I(p)=-1$.  On the other hand, if $a$ and $b$ are both
supported on $[1,\sqrt V]$, then $(a*b)(p)=0$.  Hence the exact sharp
coefficient $\mu(d)\one_I(d)$ cannot be a finite linear combination of
such balanced convolutions.  This elementary observation does not rule
out every notion of factorability, but it does rule out the balanced
shape needed for a direct invocation of the current well-factorable
distribution theorems.  Together with \eqref{eq:beta-natural-L2}, it
shows why neither a formal factor split nor coefficient $L^2$ smallness
closes the tail.
\end{remark}

The identities in this section are classical Selberg diagonalization.
The endpoint floor has close precedents in
\citet{DressIwaniecTenenbaum1983,deLaBretecheDressTenenbaum2020,RamareZunigaAlterman2026v3}.
The new role here is diagnostic: it identifies which tempting route at
the TPC-17 tail gate is unavailable before any prime-correlation
estimate is attempted.
